\chapter{검증과 결과}\label{ch:results}

이 장은 구현이 옳다는 근거와 측정 결과를 모은다. 현재는 소프트웨어 모델과 기준 RTL의 결과가 있고,
FPGA와 실물의 결과는 해당 단계가 끝난 뒤 추가한다. 이 장의 모든 수치는 저장소의 스크립트로 다시 만들 수
있으며, 명령은 부록~\ref{app:repro}에 있다.

검증의 근거는 세 가지다. 실제 전문으로 에니그마 모델을 확인하고, 무작위 키로 암호화한 전문을 봄베가
되찾는지로 봄베 모델을 확인하며, 거짓 정지의 수가 문헌의 예상과 같은 식으로 변하는지로 대각선 보드와
정지 판정을 확인한다.

\section{실제 전문}\label{sec:historical}

\paragraph{1945년 3월 6일의 M4 전문.} 독일 해군 U보트 통신망(Triton, 연합군 명칭 Shark)의 전문으로,
암호문과 키와 복호문이 공개되어 있다. 키는 얇은 반사판 B, 그리스 로터 Gamma, 로터 IV-V-III, 링
\texttt{AAPU}, 시작 위치 \texttt{ZKSA}, 스테커 10쌍이다. 에니그마 모델에 이 키를 넣어 본문 344글자를
복호화하면 공개된 복호문과 글자 단위로 일치한다. 복호문은 \texttt{EULAUFEN NACH PERIODE J BRUNO ...}로
시작하는 무선 표지 운용 시간표다. 전문의 출처와 판독상의 주의점은 \code{data/vectors/historical.json}에
적어 두었다.

이 전문을 키를 모르는 것으로 하고 봄베 방식으로 풀었다. 크립은 복호문의 48번째 글자부터 26글자
(\texttt{PEEINSEINSAMSIEBENXMAERZVO})다. 이 크립은 답을 보고 고른 것임을 밝혀 둔다. 가운데 로터의 넘어감
두 번 사이에 정확히 들어가는 구간을 골랐다. 전시의 해독자는 넘어감 위치를 알 수 없었으므로 크립을
나누거나 메뉴를 여러 벌 돌렸다(\ref{sec:turnover}절). 메뉴는 선 14개, 글자 12개, 루프 3개다. 실제 로터
순서를 포함한 여섯 개의 로터 순서(2,741,856 위치)를 돌리면 정지가 네 번 나오고 그중 하나만 체킹을
통과한다. 복원된 키는 로터 순서, 링 \texttt{AAPU}, 시작 위치 \texttt{ZKSA}, 스테커 10쌍이 모두 공개된
키와 같고 복호문도 일치한다. 전문이 길어 가운데 로터가 노치를 지나가므로 가운데 로터의 링까지 정해진다.

\paragraph{1941년 7월 8일의 육군 전문.} 북부집단군의 그날 키 표와 전문 23번의 암호문(본문 120글자)이
공개되어 있다. 공개된 평문은 찾지 못했다. 키 표의 설정으로 복호화하면 독일어 문장이 나오며, 이 결과를
모델의 회귀 시험값으로 쓴다. 로터 순서와 시작 위치는 표의 지면 배치에서 읽어 낸 것이므로 이 전문은
앞의 것보다 근거가 약하다.

\section{왕복 시험}

봄베 모델의 주된 근거는 왕복 시험이다. 임의의 키로 평문을 암호화하고, 크립과 암호문만으로 봄베를 돌려
원래 키의 로터 순서와 코어 위치가 정지 목록에 있는지, 체킹이 원래의 스테커와 평문을 되찾는지 확인한다.
\ref{ch:software}장의 예제가 그 한 경우다(표~\ref{tab:roundtrip}).

\begin{table}[ht]
\centering
\caption{예제 전문(121글자, 크립 위치 28)의 왕복 시험. 로터 순서는 Enigma~I의 120가지 전부(로터 순서 60가지 × 반사판 2종).}
\label{tab:roundtrip}
\begin{tabular}{lrrrrr}
\toprule
크립 & 메뉴의 선 & 글자 & 루프 & 정지 & 체킹 통과 \\
\midrule
\texttt{WETTERVORHERSAGE} (16글자) & 14 & 14 & 1 & 157 & 1 \\
\texttt{WETTERVORHERSAGEFUERDIE} (23글자) & 14 & 11 & 4 & 2 & 1 \\
\bottomrule
\end{tabular}
\end{table}

두 경우 모두 통과한 해는 원래의 로터 순서, 코어 위치, 스테커 10쌍, 평문 전체와 일치한다. 첫 번째 메뉴의
정지는 로터 순서당 평균 1.31회다. 이 메뉴(14글자, 루프 1개)에 대한 튜링의 추정은 약 1.6회다%
~\cite{carter-stops,wikipedia-bombe}. 스크램블러 수가 같아도 루프가 1개에서 4개로 늘면 정지가 157회에서
2회로 줄어든다. 전시의 운용 기준은 로터 순서당 정지 한 개 이상(기계가 제대로 돌았는지 확인하기 위해),
약 네 개 이하(한 사람이 체킹을 감당할 수 있는 한도)였다~\cite[31, 44, 49]{us6812report1944}.

\section{거짓 정지의 통계}\label{sec:stopstats}

정답이 없는 무작위 메뉴에서는 모든 정지가 거짓 정지다. 글자 수와 루프 수를 정한 무작위 연결 메뉴를
무작위 로터 순서 하나에 걸어 17,576 위치를 구동하고, 경우마다 400번 되풀이해 정지 수의 평균을 구했다.
그림~\ref{fig:stopstats}는 그 결과를 튜링의 추정표와 함께 그린 것이다. 추정표는 원문을 직접 열람하지
못해 2차 자료의 인용을 옮긴 값이다~\cite{carter-stops,wikipedia-bombe}.

\begin{figure}[ht]
\centering
\begin{tikzpicture}
\begin{semilogyaxis}[
    width=0.62\linewidth, height=78mm,
    xlabel={메뉴의 루프 수}, ylabel={로터 순서당 정지 수 (평균)},
    xtick={0,1,2,3}, xmin=-0.3, xmax=3.3,
    ymin=0.005, ymax=100000,
    ytick={0.01,0.1,1,10,100,1000,10000},
    yticklabels={0.01,0.1,1,10,100,1000,10000},
    grid=major, grid style={cgray!25},
    legend style={font=\footnotesize,at={(1.04,1.0)},anchor=north west,draw=cgray!50,cells={anchor=west}},
    label style={font=\small}, tick label style={font=\small},
  ]
  % 측정값 (대각선 보드 있음): data/stop_statistics.csv
  \foreach \n/\col/\mk in {8/cred/square*,10/corange/triangle*,12/cblue/*,14/cgreen/diamond*} {
    \edef\temp{\noexpand\addplot[\col,mark=\mk,line width=0.9pt,mark size=2.2pt]
      table[col sep=comma,x=loops,y=stops_db,
            restrict expr to domain={\noexpand\thisrow{letters}}{\n:\n}] {data/stop_statistics.csv};
      \noexpand\addlegendentry{측정, \n글자}}
    \temp
  }
  % 튜링의 추정: data/turing_estimate.csv
  \foreach \n/\col/\mk in {8/cred/square,10/corange/triangle,12/cblue/o,14/cgreen/diamond} {
    \edef\temp{\noexpand\addplot[\col,mark=\mk,dashed,line width=0.7pt,mark size=2.2pt,mark options={solid}]
      table[col sep=comma,x=loops,y=stops,
            restrict expr to domain={\noexpand\thisrow{letters}}{\n:\n}] {data/turing_estimate.csv};
      \noexpand\addlegendentry{튜링의 추정, \n글자}}
    \temp
  }
  \draw[cgray,dotted,line width=0.8pt] (axis cs:-0.3,17576) -- (axis cs:3.3,17576)
     node[pos=0.72,above,font=\scriptsize,text=cgray] {전체 위치 17,576};
\end{semilogyaxis}
\end{tikzpicture}
\caption{메뉴의 글자 수와 루프 수에 따른 거짓 정지 수(대각선 보드 있음). 실선은 이 프로젝트의 측정값,
점선은 튜링의 추정표다. 세로축은 로그 눈금이다.}
\label{fig:stopstats}
\end{figure}

그림에서 읽을 수 있는 것은 세 가지다.

\begin{itemize}
\item \textbf{루프 하나가 정지를 약 $1/26$로 줄인다.} 로그 눈금에서 측정값과 추정값의 기울기가 같다.
12글자 메뉴에서 루프 0, 1, 2, 3개일 때 측정값은 673, 25.3, 0.88, 0.03이고 추정값은 820, 31, 1.2, 0.04다.
이웃한 값의 비는 양쪽 모두 26 안팎이다. \ref{sec:menu}절에서 루프 하나가 ``한 바퀴 돌아 같은 선으로
돌아와야 한다''는 조건 하나에 해당하고 그것이 우연히 맞을 확률이 $1/26$이라고 한 것과 일치한다.
\item \textbf{절대값은 추정의 약 0.8배다.} 추정표가 어떤 가정으로 계산되었는지 원문으로 확인하지 못했으므로
이 차이의 원인은 단정하지 않는다. 구현 오류를 잡는 수단으로는 비율의 일치가 더 중요하다. 대각선 보드의
배선이 한 곳이라도 틀리면 루프당 감소율이 달라진다.
\item \textbf{8글자, 루프 0개에서는 비교가 성립하지 않는다.} 측정값은 위치의 수 17,576을 넘을 수 없는데
거의 모든 위치에서 정지하고 있고(17,221), 추정표의 값(40,000)은 위치의 수보다 크다. 추정표가 무엇을
세었는지 확인하지 못했다.
\end{itemize}

같은 실험을 대각선 보드 없이 하면 웰치먼의 개량이 무엇을 바꾸었는지가 보인다(그림~\ref{fig:nodiag}).
대각선 보드가 없으면 루프가 0개이거나 1개인 메뉴는 거의 모든 위치에서 정지해 기계가 아무것도 걸러 내지
못하고, 글자 수를 늘려도 나아지지 않는다. 루프가 세 개는 있어야 정지가 수십 개로 줄어든다. 대각선 보드가
없으면 루프가 세 개 이상 필요했다는 서술~\cite{carter2010bombe}과 맞는다.

\begin{figure}[ht]
\centering
\begin{tikzpicture}
\begin{semilogyaxis}[
    ybar, width=0.8\linewidth, height=62mm, bar width=7mm,
    xlabel={메뉴의 루프 수 (12글자 메뉴)}, ylabel={로터 순서당 정지 수},
    xtick={0,1,2,3}, xmin=-0.6, xmax=3.6,
    ymin=0.01, ymax=200000,
    ytick={0.01,0.1,1,10,100,1000,10000},
    yticklabels={0.01,0.1,1,10,100,1000,10000},
    ymajorgrids, grid style={cgray!25},
    legend style={font=\footnotesize,at={(0.98,0.97)},anchor=north east,draw=cgray!50},
    label style={font=\small}, tick label style={font=\small},
    nodes near coords={\pgfmathprintnumber[fixed,precision=2,1000 sep={,}]{\pgfplotspointmeta}},
    nodes near coords style={font=\scriptsize},
    point meta=rawy,
    log origin=infty,
  ]
  \addplot[fill=cgray!45,draw=cgray] table[col sep=comma,x=loops,y=stops_nodb,
     restrict expr to domain={\thisrow{letters}}{12:12}] {data/stop_statistics.csv};
  \addlegendentry{대각선 보드 없음}
  \addplot[fill=cblue!70,draw=cblue] table[col sep=comma,x=loops,y=stops_db,
     restrict expr to domain={\thisrow{letters}}{12:12}] {data/stop_statistics.csv};
  \addlegendentry{대각선 보드 있음}
\end{semilogyaxis}
\end{tikzpicture}
\caption{대각선 보드의 효과(12글자 메뉴, 경우당 400회 평균). 세로축은 로그 눈금이다.}
\label{fig:nodiag}
\end{figure}

\section{탐색 성능}\label{sec:results-search}

전체 탐색의 성능은 117글자 전문과 23글자 크립(선 14개, 글자 11개, 루프 4개의 메뉴, 크립 구간에 넘어감
없음)으로 측정했다(표~\ref{tab:perf}). 환경은 논리 스레드 24개의 PC다.

\begin{table}[ht]
\centering
\caption{소프트웨어 봄베의 전체 탐색 성능.}
\label{tab:perf}
\begin{tabular}{lrrrrr}
\toprule
기종 & 로터 순서 & 탐색 위치 & 정지 & 해 & 시간 \\
\midrule
Enigma I (I--V, 반사판 B/C) & 120 & 2,109,120 & 2 & 1 & 0.21--0.23초 \\
M4 (I--VIII, Beta/Gamma, 얇은 B/C) & 1,344 & 614,175,744 & 176 & 1 & 98초 \\
\bottomrule
\end{tabular}
\end{table}

속도는 Enigma~I에서 초당 약 900만에서 1,000만 위치, M4에서 초당 약 630만 위치다. 두 경우 모두 로터 순서,
코어 위치, 스테커 10쌍, 전문 전체가 정확히 복원되었다. 링 설정은 정해지지 않는 링을 \texttt{A}로 둔 등가
키로 보고된다(\ref{sec:sw-check}절).

비교를 위해 적으면, 실제 영국 3로터 봄베는 로터 순서 하나에 약 10.4분(정지가 없을 때)이
걸렸고~\cite{jackson-bh09,jackson-bh11}, 미 해군 4로터 봄베는 4휠 구동 1회에 약 20분이 걸렸다고
전한다~\cite{ekhall-hallenberg2018}. 뒤의 값은 직접 열람하지 못한 문헌을 거친 간접 인용이다.

\section{기준 RTL의 검증 결과}

기준 RTL은 표~\ref{tab:rtl-tests}의 시험을 모두 통과한다. 스크램블러의 치환표, 한 위치의 676선 상태 전체,
전체 스윕의 정지 목록(순서와 시험한 위치 수 포함)이 Rust 코어와 일치한다. 넘어감 가설, 대각선 보드를 끈
경우, M4가 포함된다. 따라서 ``클럭마다 한 걸음씩 퍼뜨리는 회로가 소프트웨어의 작업 목록 방식과 같은
고정점에 도달한다''는 \ref{sec:why-clock}절의 주장은 시험으로도 확인되었다.

\begin{todobox}{FPGA의 결과}
보드에서 측정한 클럭 주파수, 엔진 수, 자원 사용량, 로터 순서 하나의 스윕 시간과 전체 탐색 시간.
소프트웨어(표~\ref{tab:perf})와 실제 봄베의 구동 시간과 함께 한 표로 비교한다. 배치와 배선까지의 결과는
\ref{sec:fpga-results}절에, 세 기계의 설계값은 표~\ref{tab:hw-three}에 있다.
\end{todobox}

\begin{todobox}{실물의 결과}
기어식과 릴레이식 스크램블러의 치환 대조, 소규모 뱅크에서 같은 메뉴를 두 방식으로 돌린 정지 위치의 비교.
\end{todobox}
